% GENERATED FILE. Edit canonical lessons, then run npm run generate:books.
% canonical-source-hash: fnv1a64:4e324af73516e1bd
% canonical-lessons: IT-C463-scoglio, IT-C463-abbronzarsi, IT-C463-conchiglia, IT-C463-marea, IT-C463-tuffarsi

\chapter{Sugli scogli}
\label{ch:it-sugli-scogli}
\hlchaptermodality{463}

\begin{chapteropening}
\textbf{By the end of this chapter:} \emph{I can describe a day on a rocky shore: the rocks, getting a tan, seashells, the tide and diving in.}

Complete the last lesson of chapter 463: I can describe a day on a rocky shore: the rocks, getting a tan, seashells, the tide and diving in.
\end{chapteropening}

\section[lo scoglio]{lo scoglio — the rock (in the sea)}

\label{lesson:IT-C463-scoglio}



\begin{quote}
Before the new one: say the Italian for the cave, then the Italian for to climb.
\end{quote}

\subsection*{You'll want to know: lo scoglio}
\textbf{lo scoglio} — \textquotedblleft{}the rock (in the sea)\textquotedblright{}.

\begin{quote}
Prendiamo il sole sugli scogli.
\end{quote}

\begin{quote}
\emph{We sunbathe on the rocks.}
\end{quote}

Figuratively, an obstacle.

\begin{cousinweb}
\emph{scoglio} is inherited from Vulgar Latin \emph{scoclus}, from Latin \emph{scopulus} 'rock, crag' (EtymDB).
\end{cousinweb}

\begin{grammarlens}[title={What you've built}]
One more word for telling, explaining and describing at B1.
\end{grammarlens}

\subsection*{Your turn}
\begin{itemize}
\raggedright
  \item \emph{Say it:} \emph{lo scoglio}
  \item \emph{Say it:} \emph{lo scoglio}, once more
  \item \emph{Say it:} a sentence of your own with \emph{lo scoglio}
  \item \emph{Recall:} say the Italian for the cave, then the Italian for to climb, then say \emph{lo scoglio} again
\end{itemize}

\begin{tcolorbox}[breakable,colback=teal!4,colframe=teal!35!black,title={Before you move on}]
What does lo scoglio mean? (\textquotedblleft{}The rock (in the sea)\textquotedblright{}.) Say it once more.
\end{tcolorbox}
\section[abbronzarsi]{abbronzarsi — to get a tan}

\label{lesson:IT-C463-abbronzarsi}



\begin{quote}
Before the new one: say the Italian for to climb, then the Italian for the rock.
\end{quote}

\subsection*{You'll want to know: abbronzarsi}
\textbf{abbronzarsi} — \textquotedblleft{}to get a tan\textquotedblright{}.

\begin{quote}
Mi abbronzo subito al mare.
\end{quote}

\begin{quote}
\emph{I tan quickly at the seaside.}
\end{quote}

\emph{L'abbronzatura} — the tan.

\begin{cousinweb}
\emph{abbronzarsi} is built from \emph{a-} and \emph{bronzo} 'bronze': turning the colour of bronze. The origin of \emph{bronzo} is not settled.
\end{cousinweb}

\begin{grammarlens}[title={What you've built}]
One more word for telling, explaining and describing at B1.
\end{grammarlens}

\subsection*{Your turn}
\begin{itemize}
\raggedright
  \item \emph{Say it:} \emph{abbronzarsi}
  \item \emph{Say it:} \emph{abbronzarsi}, once more
  \item \emph{Say it:} a sentence of your own with \emph{abbronzarsi}
  \item \emph{Recall:} say the Italian for to climb, then the Italian for the rock, then say \emph{abbronzarsi} again
\end{itemize}

\begin{tcolorbox}[breakable,colback=teal!4,colframe=teal!35!black,title={Before you move on}]
What does abbronzarsi mean? (\textquotedblleft{}To get a tan\textquotedblright{}.) Say it once more.
\end{tcolorbox}
\section[la conchiglia]{la conchiglia — the seashell}

\label{lesson:IT-C463-conchiglia}



\begin{quote}
Before the new one: say the Italian for the rock, then the Italian for to get a tan.
\end{quote}

\subsection*{You'll want to know: la conchiglia}
\textbf{la conchiglia} — \textquotedblleft{}the seashell\textquotedblright{}.

\begin{quote}
I bambini raccolgono conchiglie sulla spiaggia.
\end{quote}

\begin{quote}
\emph{The children pick up shells on the beach.}
\end{quote}

\emph{Le conchiglie} are also a pasta shape.

\begin{cousinweb}
\emph{conchiglia} is inherited from Latin \emph{conchylium} 'shellfish' (Lewis \& Short), from Greek \emph{konkhúlion}, a little \emph{kónkhē} 'mussel'.
\end{cousinweb}

\begin{grammarlens}[title={What you've built}]
One more word for telling, explaining and describing at B1.
\end{grammarlens}

\subsection*{Your turn}
\begin{itemize}
\raggedright
  \item \emph{Say it:} \emph{la conchiglia}
  \item \emph{Say it:} \emph{la conchiglia}, once more
  \item \emph{Say it:} a sentence of your own with \emph{la conchiglia}
  \item \emph{Recall:} say the Italian for the rock, then the Italian for to get a tan, then say \emph{la conchiglia} again
\end{itemize}

\begin{tcolorbox}[breakable,colback=teal!4,colframe=teal!35!black,title={Before you move on}]
What does la conchiglia mean? (\textquotedblleft{}The seashell\textquotedblright{}.) Say it once more.
\end{tcolorbox}
\section[la marea]{la marea — the tide}

\label{lesson:IT-C463-marea}



\begin{quote}
Before the new one: say the Italian for to get a tan, then the Italian for the seashell.
\end{quote}

\subsection*{You'll want to know: la marea}
\textbf{la marea} — \textquotedblleft{}the tide\textquotedblright{}.

\begin{quote}
Con l'alta marea la spiaggia sparisce.
\end{quote}

\begin{quote}
\emph{At high tide the beach disappears.}
\end{quote}

\emph{L'alta marea}, \emph{la bassa marea}; also a flood of people: \emph{una marea di gente}.

\begin{cousinweb}
\emph{marea} is built on \emph{mare} 'sea', which is inherited from Latin \emph{mare} (EtymDB).
\end{cousinweb}

\begin{grammarlens}[title={What you've built}]
One more word for telling, explaining and describing at B1.
\end{grammarlens}

\subsection*{Your turn}
\begin{itemize}
\raggedright
  \item \emph{Say it:} \emph{la marea}
  \item \emph{Say it:} \emph{la marea}, once more
  \item \emph{Say it:} a sentence of your own with \emph{la marea}
  \item \emph{Recall:} say the Italian for to get a tan, then the Italian for the seashell, then say \emph{la marea} again
\end{itemize}

\begin{tcolorbox}[breakable,colback=teal!4,colframe=teal!35!black,title={Before you move on}]
What does la marea mean? (\textquotedblleft{}The tide\textquotedblright{}.) Say it once more.
\end{tcolorbox}
\section[tuffarsi]{tuffarsi — to dive, to plunge}

\label{lesson:IT-C463-tuffarsi}



\begin{quote}
Before the new one: say the Italian for the seashell, then the Italian for the tide.
\end{quote}

\subsection*{You'll want to know: tuffarsi}
\textbf{tuffarsi} — \textquotedblleft{}to dive, to plunge\textquotedblright{}.

\begin{quote}
Ci tuffiamo dallo scoglio più alto.
\end{quote}

\begin{quote}
\emph{We dive off the highest rock.}
\end{quote}

\emph{Tuffarsi in un libro} — to lose yourself in a book.

\begin{cousinweb}
\emph{tuffarsi} is the reflexive of \emph{tuffare} 'to plunge', which EtymDB traces to Lombardic \emph{taufan} 'to dip, to immerse', a Germanic word borrowed into Italian.
\end{cousinweb}

\begin{grammarlens}[title={What you've built}]
One more word for telling, explaining and describing at B1.
\end{grammarlens}

\subsection*{Your turn}
\begin{itemize}
\raggedright
  \item \emph{Say it:} \emph{tuffarsi}
  \item \emph{Say it:} \emph{tuffarsi}, once more
  \item \emph{Say it:} a sentence of your own with \emph{tuffarsi}
  \item \emph{Recall:} say the Italian for the seashell, then the Italian for the tide, then say \emph{tuffarsi} again
\end{itemize}

\begin{tcolorbox}[breakable,colback=teal!4,colframe=teal!35!black,title={Before you move on}]
What does tuffarsi mean? (\textquotedblleft{}To dive, to plunge\textquotedblright{}.) Say it once more.
\end{tcolorbox}
